\section{Structures}\label{sec:structures}

In this section, we briefly recall some basic notions from model theory and describe the main examples of structures that we will consider in this paper.

Let us begin by fixing some terminology.
A \emph{vocabulary} is a set of relations, each with a specified arity; we do not use functions in this paper. 
For example, the vocabulary of graphs contains one binary relation: the edge relation. 
(We do not include equality in the vocabulary, since it is automatically present.) 
A \emph{structure} $\A$ over a vocabulary consists of an underlying set, also denoted by $\A$, together with interpretations of relations from the vocabulary as actual relations on that set. 
An \emph{embedding} between two structures over the same vocabulary is an injective function between their underlying sets that preserves and reflects all relations. 
An \emph{isomorphism} is a surjective embedding.
An \emph{automorphism} of a structure is an isomorphism from the structure to itself; these form a group. 

Automorphisms of $\A$ act on tuples in $\A^d$ componentwise. 
When we speak of orbits in $\A^d$, we mean the orbits under this action of the automorphism group. 
For example, if $\A$ is a graph, then two pairs $(a_1,a_2)$ and $(b_1,b_2)$ are in the same orbit if and only if there is some automorphism of the graph that maps $a_1$ to $b_1$ and $a_2$ to $b_2$. 
In particular, the edge relation must be defined in the same way for both pairs;
that is, $\{a_1, a_2\} \to \{b_1, b_2\}; a_1 \mapsto b_1, a_2 \mapsto b_2$ must be an isomorphism.

All structures considered in this paper will be countable (or finite), and we will always want them to have finitely many orbits in every finite dimension as per the following definition.

\begin{definition}[Oligomorphic structure]\label{def:oligomorphic-structure}
    A structure $\A$ is \emph{oligomorphic} if $\A^d$ has finitely many orbits for every $d \in \set{1,2,\ldots}$.
\end{definition}

A relation on a structure $\A$~---~i.e.~a subset of $\A^d$~---~is called \emph{equivariant} if it is invariant under the action of the automorphism group. 
Equivalently, the relation is a union of orbits. 
If the structure is oligomorphic, then there are finitely many orbits to consider once the dimension $d$ is fixed, 
and therefore only finitely many equivariant relations. 
By the Ryll-Nardzewski Theorem~\cite[Thm.~7.3.1]{hodges1993model}, if the structure is oligomorphic and countable, 
then the equivariant relations (i.e.~subsets) are exactly those that can be defined in first-order logic~---~see for instance~\cite[Cl.~5.9]{bojanczyk_slightly}. 
In fact, the infinite structures that we consider in this paper will satisfy a stronger property: 
an equivariant relation will be definable not only by a first-order formula, but even by a quantifier-free one. 
This will be ensured by the additional homogeneity condition defined below. 
In the condition, a \emph{substructure} of $\A$ is any structure obtained by restricting $\A$ to some subset  of its underlying set; 
we do not distinguish between substructures and subsets.

\begin{definition}[Homogeneous structure] 
    A structure $\A$ is \emph{homogeneous} if every isomorphism between finite substructures of $\A$ 
    (i.e.~embedding of a finite substructure of $\A$ into $\A$) 
    extends to an automorphism of $\A$.
\end{definition}

In a homogeneous structure $\A$, an orbit in $\A^d$ consists of tuples that satisfy the same quantifier-free formulas~---~see for instance~\cite[Thm.~6.3]{bojanczyk_slightly}. 
Every homogeneous structure arises via a construction called the \emph{Fraïssé limit}~\cite[Sec.~7.4]{hodges1993model}, from classes of finite structures that satisfy certain closure properties (the so-called \emph{amalgamation classes}). 
For the Fraïssé limit to be not only homogeneous, but also oligomorphic, we need to assume that the underlying class~---~consisting precisely of the finite structures that embed in the Fraïssé limit~---~has only finitely many non-isomorphic structures of each finite size. 
This assumption is automatically satisfied if the vocabulary contains only finitely many relations of each fixed arity,
for example if the vocabulary is finite.

Here are some of the important structures that we will consider in this paper.
They are, as we will demonstrate, all homogeneous and~---~thanks to the arity-wise finiteness of the vocabularies~---~oligomorphic.
\begin{example}[Equality atoms] \label{ex:equality-atoms} 
    In this structure, the underlying set is the natural numbers and there are no relations other than equality. 
    Automorphisms are arbitrary permutations. 
    Two tuples are in the same orbit necessarily have the same equality pattern,
    and conversely this condition is sufficient~---~that is, we have homogeneity.
    For instance, given two tuples $(a_1, a_2)$ and $(b_1, b_2)$ that both have non-repeating entries, we can extend the bijection $a_1 \mapsto b_1, a_2 \mapsto b_2$ first to a permutation of $\{a_1, a_2, b_1, b_2\}$, and then to a permutation of the whole structure by mapping other elements identically.
    There is another orbit in dimension $d = 2$, defined by $x_1 = x_2$.
\lipicsEnd\end{example}

\begin{example}[Ordered atoms] \label{ex:order-atoms} 
    In this structure, the underlying set is the rational numbers, 
    equipped with the usual order~---~the vocabulary consists of this binary relation only.
    Automorphisms are order-preserving permutations. 
    Two tuples are in the same orbit if and only if they have the same order pattern. 
    Indeed, given two tuples $(a_1, a_2)$ and $(b_1, b_2)$ both with strictly increasing entries, we can extend $a_1 \mapsto b_1, a_2 \mapsto b_2$ to an automorphism by using Cantor's back-and-forth method, or by explicitly defining a monotone piecewise linear bijection.
    There are two other orbits in dimension $d = 2$, defined by $x_1 = x_2$ and $x_1 > x_2$.
\lipicsEnd\end{example}

\begin{example}[Vector atoms] \label{ex:vector-space-atoms} 
    Fix some finite field $\innerfield$, and let $\A$ be the vector space of countable dimension over $\innerfield$. 
    This vector space is seen as a structure over an infinite vocabulary containing, for every $d \in \set{1,2,\dots}$ and coefficients $\lambda_1, \dots, \lambda_d$ in $\innerfield$, the relation 
    \begin{align*}
        \setbuild{(a_1,\dots,a_d) \in \A^d}{ $\lambda_1 a_1 + \cdots + \lambda_d a_d = 0$ }.
    \end{align*}
    The vocabulary is defined so that automorphisms are the same thing as permutations that are linear maps. 
    Two tuples of the same length are in the same orbit if and only if they have the same linear dependencies. 
    By way of illustration, consider two tuples $(a_1, a_2)$ and $(b_1, b_2)$ each consisting of linearly independent entries.
    Extend $a_1, a_2$ to a basis for the subspace $V$ of $\A$ spanned by $a_1, a_2, b_1, b_2$; 
    do the same for $b_1, b_2$.
    This gives us an automorphism of $V$ with $a_1 \mapsto b_1, a_2 \mapsto b_2$, which can be then extended to an automorphism of $\A$ by mapping other vectors identically.
    In $\A^2$ there are $2 + |\innerfield|$ more orbits:
    $x_1 = 0 = x_2$;
    $x_1 = 0 \neq x_2$;
    and $0 \neq x_1 = \lambda x_2$, where $\lambda$ ranges over the field $\innerfield$.
\lipicsEnd\end{example}

\begin{example}[Rado atoms]\label{ex:rado-graph} 
    The Rado graph is the Fraïssé limit of the class of finite undirected graphs. 
    Here, an undirected graph is viewed as a structure with one binary relation that is symmetric and irreflexive. 
    One explicit description of the Rado graph is as follows: its vertices are the natural numbers, and there is an edge between $n < m$ if the $n$-th least significant bit in the binary representation of $m$ is $1$.
    But a famous characterisation of the Rado graph is that if one randomly selects a graph with a countable set of vertices by independently including each possible edge with probability $1/2$, then with probability $1$ the resulting graph is isomorphic to the Rado graph. 
    
    As a Fraïssé limit, the Rado graph is homogeneous by construction.
    So two tuples are in the same orbit if and only if they have the same equality and adjacency patterns.
    In dimension $d = 2$ there are three orbits:
    the two coordinates can be equal, adjacent hence distinct, or distinct but non-adjacent.
\lipicsEnd\end{example}
